\section*{Step 265: P-vs-NP Attack Foreclosure Replication}

\paragraph{Verdict.}
\[
\mathrm{V\_PNP\_attack\_foreclosure\_replicated}.
\]
The Attack Foreclosure Conjecture is replicated on the P-vs-NP track, raising
the candidate meta-theorem to verified-on-5-tracks:
\[
\mathrm{RH,\ BSD,\ Hodge,\ NS,\ P\mbox{-}vs\mbox{-}NP}.
\]

\paragraph{P-vs-NP Attack Foreclosure Conjecture.}
Under the Six Birds Foundations III typed-condition discipline, no proof of
\(P=NP\) or \(P\neq NP\) can be constructed by carrier pivoting, bridge import,
or cascade-internal computation alone, in isolation from named external
classical content.

\paragraph{Move (a): carrier pivoting.}
The P-vs-NP carrier atlas includes lawful declared packages, fixed quotient
non-descent, scoped local-status Tseitin obstructions, proof-system bridge
leaves, \(\Xi_{\mathrm{atlas}}(P)\), universal \(P\)-machine atlas
target-equivalence, packaging-axis level-1 hosts, abstract-transformer package
atlases \(\Pi_r^\#(\varphi)=\operatorname{lfp}(F_{\varphi,r}^\#)\),
no-smuggling/readability/atlas-scope gates, barrier stacks, and algorithmic/GCT
route-specific gates.  Pivoting among these carriers preserves target-strength,
package-atlas terminality, or bridge-failure status.

\paragraph{Move (b): bridge import.}
Cook-Levin and Karp reductions organize NP-completeness but do not decide
P-vs-NP.  Baker-Gill-Solovay relativization, Razborov-Rudich natural proofs,
and Aaronson-Wigderson algebrization are barrier records.  Williams' ACC lower
bounds and Mulmuley-Sohoni GCT are restricted or route-specific programs.
Transport from any such adjacent result to full P-vs-NP is P-vs-NP-strength.

\paragraph{Move (c): cascade-internal computation.}
The inherited package-atlas CTMT chain refines
\[
\Xi_{\mathrm{pack}}
\to \text{fixed quotient}
\to \text{Tseitin}
\to \text{proof-system leaves}
\to \Xi_{\mathrm{atlas}}(P)
\to \text{universal \(P\)-machine TE no-go}
\to \text{level-1 host}
\to \Pi_r^\#
\to \text{no-smuggling/readability/atlas-scope}
\to \text{barrier stack}
\to \text{algorithmic/GCT gates}.
\]
The computation exposes package-lawfulness, barrier, and route-specific terminal
objects rather than a proof-complete resolution of P-vs-NP.

\paragraph{P-vs-NP-specific audit.}
Diagonalization maps to cascade computation and is bounded by relativization
unless nonrelativizing content is supplied.  Reductions map to bridge import.
Circuit lower bounds, proof complexity, communication complexity, average-case
complexity, and adjacent quantum models are carrier pivots or bridge imports.
GCT and topological/homological/algebraic methods are cascade routes with
external mathematical gates.  No fourth standard internal move emerged.

\paragraph{Not foreclosed.}
The result preserves NP48 package specification, non-circular P-visibility
atlases, witness-fiber closure deficits, proof-system-to-packaging bridges,
GCT/algorithmic structural bridges, barrier-avoiding methods, fresh complexity
carriers, fresh typed mechanisms, and any actual P-vs-NP proof in either
direction.

\paragraph{Classical references.}
The audit cites Cook 1971, Karp 1972, Baker-Gill-Solovay 1975,
Razborov-Rudich 1994/1997, Aaronson-Wigderson 2008, Williams 2011/2014,
Mulmuley-Sohoni 2001, and Grigoriev/Schoenebeck proof-system records.
